% !TEX root = main.tex
% Outils de mise en page : reprise de ceux de la soutenance de stage
% (diagramoverlay, callouts, pastilles, fiche d'orateur), plus quelques
% outils propres au projet (captures, classes UML, puces de fichier).

% Hauteur utile pour une image : la page moins les deux bandeaux, moins un jeu.
\newlength{\astdiagheight}
\setlength{\astdiagheight}{\dimexpr\paperheight-\astheadheight-\astfootheight-0.7cm\relax}

% Separateur entre mots-cles :  surface \astsep fuite
% \texorpdfstring : dans les metadonnees du PDF (\author -> pdfauthor), une
% simple virgule, sinon hyperref avertit << Token not allowed in a PDF string >>
\newcommand{\astsep}{\texorpdfstring{\;\textperiodcentered\;}{, }}

% ---------------------------------------------------------------------------
% diagramoverlay : on dessine sur une image en coordonnees relatives,
% (0,0) = coin bas gauche, (1,1) = coin haut droit.
% ---------------------------------------------------------------------------
\newif\ifastdiaggrid
\newcommand{\astdrawgrid}{%
  \draw[cyan!60,very thin,step=0.1] (0,0) grid (1,1);
  \foreach \i in {0,0.1,...,1.01}{%
    \node[font=\tiny,text=cyan!70!black,anchor=north] at (\i,0)
      {\pgfmathprintnumber[fixed,precision=1]{\i}};
    \node[font=\tiny,text=cyan!70!black,anchor=east] at (0,\i)
      {\pgfmathprintnumber[fixed,precision=1]{\i}};}}

\NewDocumentEnvironment{diagramoverlay}{O{} m}{%
  \begin{center}%
  \begin{tikzpicture}
    \node[anchor=south west,inner sep=0,outer sep=0] (img) at (0,0)
      {\includegraphics[width=\textwidth,height=\astdiagheight,keepaspectratio,#1]{#2}};
    \begin{scope}[x={(img.south east)},y={(img.north west)}]
      \ifastdiaggrid\astdrawgrid\fi
}{%
    \end{scope}
  \end{tikzpicture}%
  \end{center}%
}

\tikzset{astcallout/.style={draw=phelmared,fill=white,text=inpnavy,rounded corners=1.5pt,
  line width=0.8pt,font=\footnotesize,inner sep=3pt,align=left}}
\NewDocumentCommand{\highlight}{D<>{1-} O{phelmared} m m m m}{%
  \only<#1>{\draw[#2,line width=1.2pt,rounded corners=2pt] (#3,#4) rectangle (#5,#6);}}
\NewDocumentCommand{\spotlight}{D<>{1-} m m m m}{%
  \only<#1>{\fill[white,fill opacity=0.66,even odd rule]
    (0,0) rectangle (1,1) (#2,#3) rectangle (#4,#5);}}
\tikzset{astnum/.style={circle,fill=phelmared,draw=white,line width=0.6pt,text=white,
  font=\scriptsize\bfseries,inner sep=0pt,minimum size=4.2mm}}
\newcommand{\astmark}[3]{\node[astnum] at (#1,#2) {#3};}
\NewDocumentCommand{\callout}{D<>{1-} O{} m m m}{%
  \only<#1>{\node[astcallout,#2] at (#3,#4) {#5};}}
% Etiquette reliee par une fleche :  \astpoint{x}{y}{xe}{ye}{ancre}{texte}
\newcommand{\astpoint}[6]{%
  \node[astcallout,anchor=#5] (astpt) at (#3,#4) {#6};
  \draw[-{Latex[length=1.6mm]},phelmared,line width=0.8pt] (astpt.#5) -- (#1,#2);}
\newcommand{\astpill}[1]{\tikz[baseline=(astp.base)]\node[astnum] (astp) {#1};}

% ---------------------------------------------------------------------------
% Projet : puces de code et de fichier
% ---------------------------------------------------------------------------
% \code{dynamic_cast}   mot de code dans le texte
\newcommand{\code}[1]{{\ttfamily\color{combislate}#1}}
% \filechip{world.hpp}  petite etiquette de fichier, grise
\newcommand{\filechip}[1]{\tikz[baseline=(fc.base)]\node[fill=asttrack,text=textgrey,
  rounded corners=1pt,inner xsep=2pt,inner ysep=1.2pt,font=\ttfamily\tiny] (fc) {#1};}

% ---------------------------------------------------------------------------
% Classes UML dans le langage des schemas : bord colore, fond pale.
%   \node[cls=memborder/memfill] (P) {\clsname{Particle}\clsbody{pos, speed\\mass}};
% ---------------------------------------------------------------------------
\tikzset{
  cls/.style args={#1/#2}{draw=#1,fill=#2,line width=0.8pt,rounded corners=1.6mm,
    align=center,inner sep=1.4mm,font=\scriptsize,text=astbodyfg},
  abstract/.style={dash pattern=on 3pt off 1.5pt},
  blackbox/.style={draw=floorgrey,fill=asttrack,line width=0.8pt,rounded corners=1.6mm,
    align=center,inner sep=1.4mm,font=\scriptsize,text=textgrey},
  inherit/.style={-{Triangle[open,length=2.4mm,width=2.4mm]},line width=0.7pt,draw=combislate},
  vinherit/.style={inherit,draw=phelmared,line width=1pt},
  owns/.style={{Diamond[length=2.6mm,width=1.8mm]}-{Latex[length=1.6mm]},line width=0.7pt,draw=combislate},
  uses/.style={-{Latex[length=1.6mm]},line width=0.7pt,draw=combislate,dash pattern=on 2pt off 1.5pt},
  flow/.style={-{Latex[length=2mm]},line width=1pt,draw=inpnavy},
}
\newcommand{\clsname}[1]{{\bfseries\small #1}}
\newcommand{\clsabstract}[1]{{\bfseries\itshape\small #1}}
% Le corps passe par un tabular : dans un noeud TikZ, un \\ remet la police
% du noeud, et un \\ dans un groupe casse le noeud.
\newcommand{\clsbody}[1]{\\[-0.2ex]{\color{textgrey}\tiny
  \renewcommand{\arraystretch}{0.9}\begin{tabular}{@{}c@{}}#1\end{tabular}}}
\newcommand{\clsleaf}[1]{{\bfseries\footnotesize #1}}

% ---------------------------------------------------------------------------
% \scenecard{fichier}{nom}{hauteur}  capture d'une scene dans un cadre
% d'etage pointille, avec le nom sur le bord bas (comme les etages des schemas)
% ---------------------------------------------------------------------------
\newcommand{\scenecard}[3]{%
  \begin{tikzpicture}
    \node[inner sep=0] (sc) {\includegraphics[height=#3]{#1}};
    \draw[floorgrey,line width=0.8pt,dash pattern=on 3pt off 2pt,rounded corners=1.5mm]
      ($(sc.south west)+(-1mm,-1mm)$) rectangle ($(sc.north east)+(1mm,1mm)$);
    \node[fill=white,draw=combislate,line width=0.6pt,font=\tiny\bfseries,
      inner xsep=1.5mm,inner ysep=0.8mm] at ($(sc.south)+(0,-1mm)$) {\MakeUppercase{#2}};
  \end{tikzpicture}}

% ---------------------------------------------------------------------------
% \astbig[couleur]{400}{particules}   grand nombre et sa legende
% ---------------------------------------------------------------------------
\NewDocumentCommand{\astbig}{O{inpnavy} m m}{%
  \begin{center}
    {\fontsize{24}{28}\selectfont\bfseries\color{#1}#2}\\[0.3mm]
    {\footnotesize\color{textgrey}#3}
  \end{center}}

% ---------------------------------------------------------------------------
% Fiche d'orateur (cible << fiche >>) : pour chaque diapositive, quelques
% reperes courts, a imprimer et a garder a cote.
%   \astcue{ \cue{idee, \cuekey{mot cle}} \cuewarn{piege} \cuenext{transition} }
% ---------------------------------------------------------------------------
\newcommand{\cuekey}[1]{\textbf{\color{phelmared}#1}}
\newcommand{\cue}[1]{\item[{\color{phelmared}$\rightarrow$}] #1}
\newcommand{\cuewarn}[1]{\item[{\color{dummyorange}\footnotesize\faExclamationTriangle}] {\color{dummyorange!80!black}#1}}
\newcommand{\cuenext}[1]{\item[{\color{astblue}\footnotesize\faShare}] {\color{astblue}\itshape #1}}
\ifdefined\astfichemode
  \setbeamertemplate{note page}{%
    \vspace*{2mm}%
    \begin{minipage}[t]{0.36\linewidth}\vspace{0pt}%
      \fbox{\insertslideintonotes{0.30}}%
    \end{minipage}\hspace{3mm}%
    \begin{minipage}[t]{0.58\linewidth}\vspace{1mm}%
      {\normalsize\bfseries\color{inpnavy}\insertframetitle}\par\vspace{1mm}%
      {\scriptsize\color{textgrey}diapositive \insertframenumber}%
    \end{minipage}\par\vspace{2mm}%
    \begin{minipage}[t]{0.97\linewidth}%
      \small\raggedright\insertnote
    \end{minipage}}
  \newcommand{\astcue}[1]{\note{\begin{itemize}\setlength{\itemsep}{1.2mm}#1\end{itemize}}}
\else
  \newcommand{\astcue}[1]{}
\fi

% ---------------------------------------------------------------------------
% Fil rouge de la partie 2 : les huit etapes d'un pas de simulation, en bas
% de la diapositive. L'etape courante est pleine, les suivantes estompees.
%   \fil{4}   a placer dans la frame (dessin en surimpression)
% ---------------------------------------------------------------------------
\newcommand{\filsteps}{Scène,Particules,World,Forces,Intégration,Contraintes,Dessin,Souris}
\newcommand{\fil}[1]{%
  \begin{tikzpicture}[remember picture,overlay]
    \begin{scope}[shift={($(current page.south)+(0,\astfootheight+1.6mm)$)}]
      \foreach \lab [count=\k] in \filsteps {
        \pgfmathsetmacro{\fx}{(\k-4.5)*1.72}
        \ifnum\k<#1\relax
          \node[draw=inpnavy,fill=white,text=inpnavy,rounded corners=1pt,font=\tiny,
            minimum width=1.55cm,inner ysep=1pt] at (\fx,0) {\k\enspace\lab};
        \else\ifnum\k=#1\relax
          \node[draw=phelmared,fill=phelmared,text=white,rounded corners=1pt,font=\tiny\bfseries,
            minimum width=1.55cm,inner ysep=1pt] at (\fx,0) {\k\enspace\lab};
        \else
          \node[draw=astmutedbadge,fill=white,text=floorgrey,rounded corners=1pt,font=\tiny,
            minimum width=1.55cm,inner ysep=1pt] at (\fx,0) {\k\enspace\lab};
        \fi\fi
      }
    \end{scope}
  \end{tikzpicture}}

% ---------------------------------------------------------------------------
% Schemas draw.io : sources dans schemas/*.drawio, PDF exportes par
% << make schemas >> dans figures/schemas/. Ils sont dessines a l'echelle 2 :
% inclus a 0,5, leurs textes ont la taille des textes des diapositives
% (21 px = \scriptsize). Ils ne sont reduits que s'ils depassent la place.
%   \schema{hierarchie}          hauteur max : toute la zone utile
%   \schema[5.3cm]{forces}       hauteur max donnee (diapos avec fil rouge)
% ---------------------------------------------------------------------------
\NewDocumentCommand{\schema}{O{\astdiagheight} m}{%
  \adjustbox{scale=0.5,max width=\linewidth,max totalheight=#1}%
    {\includegraphics{figures/schemas/#2.pdf}}}
